% Proposed accessible fragment; use the host edition's sans-serif text/math setup.
% Registry revision 20; SHA-256 cc9676276550560eddc941ee08874019d7566486be709a439c732b0996a1c1e8.
\section{Exactly What the LINK Controls Change}

A model can improve its average by doing better on one recording and worse
on another. That is why the LINK comparison keeps four results visible:
two finger groups on each of two validation dates. A better average cannot
cancel a harmed group.

The study pairs recorded neural features with measured positions of the index
finger and the combined middle, ring and small-finger group in one macaque.
It trains on January 27 and 30 and February 4 and 5, 2020, then selects models
on February 6 and 11. Those positions retain the source's normalized units.
They are not forces or a reconstruction of every finger joint. There is no
simultaneous muscle, peripheral-nerve or spinal recording in this comparison.

Every example contains eleven successive neural bins, ending at the position
being predicted. A twelfth bin supplies the next-neural target used by the
extra-task model. All twelve must be in the same trial, with adjacent recorded
times between 18.5 and 21.5 milliseconds apart. Even a baseline that never
uses the twelfth bin needs it to be present. Everyone is therefore compared
on the same examples. The ten earlier gaps span a history before the current
bin; the nominal 20-millisecond interval does not turn that history into a
measured nerve-conduction delay.

The 192 neural inputs contain 96 transformed threshold-crossing counts and
96 archived power features. Both sets follow the same physical electrode
contact order. A contact is not a permanently identified individual neuron.
Scaling uses only the four training dates, counting each trial-owned frame
once. Repeated overlapping windows do not give that frame extra influence
over the scale, and validation answers do not supply scaling statistics.

The first training date fixes the contact order. The count and power columns
move together when reordered, so both descriptions still refer to the same
declared physical contacts. The count transformation takes the natural
logarithm of one plus each nonnegative count. It keeps a zero count at zero
while compressing larger counts. The archived power features keep their
supplied numerical scale; the two feature types are not converted into one
physical quantity merely because they enter the same model.

The scaler calculates a separate mean and population standard deviation for
each transformed feature. Population standard deviation means subtracting the
mean, squaring those differences, averaging over all the fitting frames and
taking the square root. Those calculations use float64, the declared numerical
format with more precision than float32. Each unique frame belonging to an
eligible training trial contributes once across the training dates, even when
many overlapping windows later reuse it. The scaler therefore describes those
frames, rather than a new weighting produced by how often they appear in
windows. Validation supplies neither a mean nor a spread.

Normalization subtracts the training mean and divides by the saved training
spread. Exactly zero spread is replaced by a divisor of one, which avoids
division by zero for a constant training feature. There is no other minimum
spread: a small nonzero value is not raised to a chosen floor. There is no
clipping to trim large normalized values. This is the exact retained numerical
rule, not a general recommendation to change it after inspecting a result.
Replacing a zero divisor does not fill a missing observation or prove that
the underlying biological signal is constant.

The transform is calculated in float64, then its output is cast to float32
for the model. Casting means representing the computed values in the other
numerical format; rounding can occur. It does not improve an electrode's
physical accuracy or settle an uncertain unit. The position targets keep
their archived float32 values. The constant reference uses the mean of unique
trial-owned training positions. Nonfinite measurements, unknown contacts or
incompatible data descriptions cause a failure, rather than silently dropping
rows or filling them with invented values.

The source descriptions leave several conflicts visible. A threshold-crossing
attribute says volts, while the quantity is interpreted as event counts; the
comparison makes no voltage conversion. It also retains the archived power
values without an unexplained multiplication by four in the author's loader.
Neither choice establishes absolute power calibration. The documented
32-millisecond interval differs from the retained nominal 20-millisecond
clock. The observed timestamps and their allowed jitter define this recipe;
the conflict remains unresolved, rather than being repaired by silently
changing the interval. Full author trial counts must agree with trial-stop
timestamps, and time-series references that are uniformly one entry short
are identified separately. Changing units or discarding an inconvenient
boundary would change which experiment was being compared.


The recurrent model reads each window from its declared starting state.
Its numerical memory is reset between windows; it does not carry the previous
window's final state into the next one. Adjacent windows may reuse some source
samples, but that overlap is not an extra memory channel. Each prediction
must be accounted for by the input window and the fixed model it declares.

Within the window, each new bin updates the hidden values. The position output
combines the final hidden values with another description: for every neural
feature, its mean across all eleven input bins. That mean retains the feature's
average level while discarding its within-window order. The final recurrent
state can depend on the sequence. The output therefore has access to both
descriptions; a benefit from history alone cannot be attributed uniquely to
recurrent ordering. That is one reason for the separate order control.

A separate linear output predicts the next neural feature vector. It remains
present in both candidate architectures, but has zero objective weight for
position-only fitting. Its presence does not give that model an additional
movement target. Likewise, an old software name for the position output does
not turn the normalized position estimate into velocity. These statements
describe the retained recipe; they do not independently verify every state
reset, prediction or update in the completed execution.

The first control asks what happens when the training pairing is broken.
Imagine each example's neural history and its next-neural target held together
as one packet. The control moves that packet to another example within the
same training date, while the position answer stays where it was. It moves
whole eligible windows, not individual channels or bins, and not whole trials
as blocks. Overlapping windows can move separately. The same fixed shuffle
serves every epoch of that run. It is shared by the two model types starting
from the same seed, but differs across seeds and recording dates.

The two validation recordings remain unchanged for this control. We therefore
fit on mismatched training pairs and ask how the resulting model predicts the
ordinary measured validation positions. Some windows can accidentally stay
in place; common task structure also remains. The control challenges the
declared association without proving that every other association has gone.

The second control asks whether the declared order helps. Number an input's
bins from zero to ten, with ten being the current bin. The control presents
them as nine, eight, seven, six, five, four, three, two, one, zero, ten.
The ten older bins reverse; the current bin stays last. Position answers and
next-neural targets remain unchanged. This happens in both training and
validation. The operation is fixed, uses no random draw and preserves all
the values in each window. A model could learn the reversed arrangement,
so this is a particular order comparison rather than a guarantee that timing
information has vanished.

Both model types start from each of three fixed seeds: 20261007, 20261008 and
20261009. The starting tensors are checked against the saved initialization
for that seed. A separate generator restarts from the same seed in each
paired run and draws one training-row order per epoch. For the association
shuffle, another generator starts at the seed plus one million and advances
across the four dates. Matching seeds specify matching procedures; they do
not themselves prove that the completed execution followed every procedure.

There are 40 epochs per run, with batches of 256 and a retained final partial
batch. Every epoch remains in the record. Each seed selects its lowest
average validation error across the four date/group cells, taking the earlier
epoch on an exact tie. Final predictions average all three selected seed
models equally. Neither control competes to become the chosen measured model.
Two model types, three seeds and three conditions produce the reported
18 runs and 720 saved epochs.

The simple linear references try five fixed penalties: 0.01, 0.1, 1, 10 and
100. A tie selects the smaller penalty, then the current-bin reference ahead
of the history reference if their selected errors tie. A tie between the
two measured recurrent ensembles selects position-only. These rules prevent
a favourable seed, epoch or tie decision being quietly substituted later.

The protocol asks five further questions: does the position-only model beat
current-bin ridge; does it beat its reversed-history control; does the
extra-task model beat current-bin ridge; does it beat its own reversed-history
control; and does the extra-task model beat position-only? Each
comparison requires all three conditions together: a strictly lower average
error, no higher error in any date/group, and a negative upper end of its
paired 95-percent error-difference interval. Lower error is better, so a
candidate-minus-reference interval wholly below zero favours the candidate.

That interval repeatedly samples whole supported trials within each validation
date, keeping every sampled trial's supported frames and using the same draws
for the models and groups. It uses 1,000 draws with fixed seed 20261007 and
linear percentile interpolation. It describes uncertainty within these
validation-selected comparisons; it does not confirm a result in a new animal
or on a reserved recording.

For model admission, the three-part comparison is made against the selected
best linear reference. Every date/group must also have a defined, positive
ordinary prediction score, called $R^2$, and exceed its matched association
control. Ordinary $R^2$ compares error with predicting
that validation group's mean; the separately named training-mean score uses
a different reference. An undefined value cannot count as a pass. This
model-admission rule does not replace the five separate benefit comparisons.

The retained independent postfit checker records each condition separately
for all five comparisons. Every comparison has a lower average error, yet
every one makes at least one date/group worse. Three intervals have an upper
end below zero. None of that cancels the failed no-harm conditions: all five
benefit claims fail, and both models fail their separate admission rules.
These are the existing decisions, not a newly calculated result.

The remaining checks concern how the learning calculation was performed. The
order-control ledgers have empty provenance fields; the source and earlier
control evidence do not provide a new dynamic reversal check. The correctly
negative postfit report also cannot independently verify every optimizer
update, checkpoint or restored random state. A completed run and a saved file
are different evidence from that verification.

The later recordings remain closed: February 22 and 24, 2020, and June 21 and
22, 2023. Here, ``earlier input'' means earlier in the supplied file's source
indices. It does not establish that archive filtering used no later signal,
or that an online biological loop has the same delay. Predicting recorded
positions is a bounded decoding test. It demonstrates neither sensory
delivery nor a nervous system learning and retaining a skill.
